\begin{solucion}
    Repitiendo nuevamente el procedimiento del literal anterior:
\begin{itemize}
    \item \textbf{Polinomio mínimo de $\sqrt{3}$}: Para $\gamma = \sqrt{3}$, el polinomio $m_\gamma(x) = x^2 - 3$ es irreducible sobre $\mathbb{Q}$, ya que $3$ no es un cuadrado perfecto en $\mathbb{Q}$, o dicho de otro modo $x^2 - 3 = 0$ no tiene solución racional. Por tanto, $[\mathbb{Q}(\sqrt{3}) : \mathbb{Q}] = 2$.

    \item \textbf{Polinomio mínimo de $\sqrt[3]{2}$}: Sea $\delta = \sqrt[3]{2}$ (la raíz cúbica real de $2$). Su polinomio mínimo es $m_\delta(x) = x^3 - 2$, el cual es irreducible en $\mathbb{Q}[x]$ por el criterio de Eisenstein con $p=2$, pues $2\not|=a_n=1$ y $a|a_0=2, a|a_1=0, a|a_2=0$ y $2^2=4\not|=a_0=2$. Por lo tanto, $[\mathbb{Q}(\sqrt[3]{2}) : \mathbb{Q}] = 3$.
\end{itemize}

Dada entonces ahora la extensión $\mathbb{Q}(\sqrt{3}, \sqrt[3]{2})$. Observemos que ninguno de los dos elementos de la extension de uno está contenido en el campo generado por la extension del otro, esto es: $\sqrt{3} \notin \mathbb{Q}(\sqrt[3]{2})$ y $\sqrt[3]{2} \notin \mathbb{Q}(\sqrt{3})$. Esto es debido a que $\mathbb{Q}(\sqrt{3})$ tiene grado $2$ y $\mathbb{Q}(\sqrt[3]{2})$ tiene grado $3$ sobre $\mathbb{Q}$. No puede existir una inclusión estricta entre ellas porque, pues en caso contrario de existir $\mathbb Q \subset \mathbb{Q}(\sqrt{3}) \subset \mathbb{Q}(\sqrt[3]{2})$, entonces por el \textbf{teorema de la torre} tendríamos que 
\begin{align*}
    [\mathbb{Q}(\sqrt[3]{2}) : \mathbb{Q}] &=3= [\mathbb{Q}(\sqrt[3]{2}) : \mathbb{Q}(\sqrt{3})] \cdot [\mathbb{Q}(\sqrt{3}) : \mathbb{Q}]\\
    &= [\mathbb{Q}(\sqrt[3]{2}) : \mathbb{Q}(\sqrt{3})] \cdot 2\\
\end{align*}

Por lo que el grado $2$ tendría que dividir al grado $3$,contradicción. En consecuencia,  
 $\mathbb{Q}(\sqrt{3})\cap \mathbb{Q}(\sqrt[3]{2}) = \mathbb{Q}$, Es decir, el polinomio $x^3 - 2$ no adquiere raíces en $\mathbb{Q}(\sqrt{3})$ (y $x^2 - 3$ no adquiere raíces en $\mathbb{Q}(\sqrt[3]{2})$).

Aplicamos el \textbf{Teorema de la Torre} con la cadena:
\[
\mathbb{Q} \subset \mathbb{Q}(\sqrt{3}) \subset \mathbb{Q}(\sqrt{3},\, \sqrt[3]{2})~.
\]
Tenemos $[\mathbb{Q}(\sqrt{3}) : \mathbb{Q}] = 2$. Y como ademas:
\[
[\mathbb{Q}(\sqrt{3}, \sqrt[3]{2}) : \mathbb{Q}(\sqrt{3})]
\]
es igual al grado del polinomio mínimo de $\sqrt[3]{2}$ sobre $\mathbb{Q}(\sqrt{3})$.Pues 
 como mencionamos anteriormente, $x^3 - 2$ sigue siendo irreducible sobre $\mathbb{Q}(\sqrt{3})$, por lo que este grado es $3$. Por lo tanto:

 \begin{respuesta}
    \[
[\mathbb{Q}(\sqrt{3}, \sqrt[3]{2}) : \mathbb{Q}] = 2 \cdot 3 = 6.
\]
 \end{respuesta}


% En conclusión, el grado de la extensión $\mathbb{Q}(\sqrt{3}, \sqrt[3]{2})$ sobre $\mathbb{Q}$ es $6$.

% Una base de $\mathbb{Q}(\sqrt{3}, \sqrt[3]{2})$ sobre $\mathbb{Q}$ la podemos tomar como el producto cartesiano:
% \[
% \left\{1,\ \sqrt{3}\right\} \times \left\{1,\ \sqrt[3]{2},\ (\sqrt[3]{2})^2\right\},
% \]
% es decir:
% \[
% \left\{1,\ \sqrt{3},\ \sqrt[3]{2},\ \sqrt{3}\sqrt[3]{2},\ (\sqrt[3]{2})^2,\ \sqrt{3}(\sqrt[3]{2})^2\right\}.
% \]
% La independencia lineal de estos elementos se justifica porque provienen de combinar bases independientes de cada subextensión. Esto confirma que el grado es $6$.
\end{solucion}
